\documentclass[11pt]{article}
\usepackage{amsmath,amssymb,amsthm}
\newtheorem{theorem}{Theorem}
\newtheorem{proposition}{Proposition}
\newtheorem{lemma}{Lemma}
\begin{document}
\title{Odd targets and positive tangent kernels}
\author{}
\date{}
\maketitle

\paragraph{Status and scope.}
The forward normalization for all original representations of odd Boolean
targets is checked in \texttt{Results/OddKernelNormalization.lean}, with
main declaration \texttt{HeadComplexity.OddKernel.of\_fracComputable}.
The exact atom identity is \texttt{atom\_odd\_identity}, positivity of
the kernel denominator is \texttt{value\_den\_pos}, and strictification
under oddness is \texttt{strict\_odd\_part}, all in the same namespace.
Their axiom audits report only \texttt{propext}, \texttt{Classical.choice},
and \texttt{Quot.sound}. Mixed atom orientations and constant denominators
are included. The converse and the Fourier statements below are complete
written proofs, with additional exact rational checks, but are not formalized.
No improved head-count lower bound or upper bound follows here. In particular, the
positivity statements below do not control arbitrary signed numerator
directions or cancellations between heads.

Write $t_i=2x_i-1\in\{-1,1\}$, and let $-t$ denote the simultaneous
complement of every input bit. A target $f:\{-1,1\}^n\to\{-1,1\}$ is
odd if $f(-t)=-f(t)$. A score represents it in the original convention
when $F(t)>0$ exactly at its positive labels, and $F(t)\leq0$ at its
negative labels.

\begin{theorem}[Global positive-kernel normalization]
Suppose an odd target is represented by a bias plus $H$ original atoms.
There exist vectors $u_h\in\mathbb R^n$ and $w_h\in[0,\infty)^n$ such that
\[
 \sum_i w_{hi}<1,\qquad
 G(t)=\sum_{h=1}^H\frac{u_h\cdot t}{1-(w_h\cdot t)^2},\qquad
 f(t)G(t)>0\quad\hbox{for every }t\in\{-1,1\}^n.
\]
For each $h$, either every coordinate of $w_h$ is strictly positive,
or $w_h=0$. No strictification of the original representation is needed.
Conversely, every finite strictly sign-representing score of the displayed
form, allowing arbitrary weakly positive $w_h$ with the stated strict
sum bound, is represented by at most $2H$ original atoms. This converse
preserves its strict signs, rather than necessarily its exact values.
\end{theorem}

\begin{proof}
For one original atom, abbreviate its parameters by
$\eta,\delta,\gamma,\alpha,\rho_i,m_i$, with
$\gamma,\alpha,\rho_i>0$. Its raw-bit affine coefficients are
\[
 B(x)=\gamma+\sum_i\rho_i+\sum_i\rho_i(\alpha-1)x_i,
 \qquad
 A(x)=\eta+\sum_i\rho_i m_i+
       \sum_i\rho_i((\alpha-1)m_i+\alpha\delta)x_i.
\]
Put $R=\sum_i\rho_i$ and
\[
 d=\gamma+\frac{\alpha+1}{2}R>0,\qquad
 v_i=\frac{\rho_i(\alpha-1)}{2d},\qquad
 a_i=\frac{\rho_i((\alpha-1)m_i+\alpha\delta)}{2d},
 \qquad
 a_0=\frac{\eta+\sum_i\rho_i m_i}{d}+\sum_i a_i.
\]
Then the atom is exactly
\[
 \frac{a_0+a\cdot t}{1+v\cdot t}.
\]
All $v_i$ have the same strict sign when $\alpha\ne1$, and $v=0$ when
$\alpha=1$. Moreover
\[
 1-\sum_i|v_i|
 =\frac{2\gamma+(\alpha+1-|\alpha-1|)R}{2d}>0.
\]
The last numerator is positive since
$\alpha+1-|\alpha-1|=2\min(\alpha,1)>0$.
In particular $1\pm v\cdot t>0$ on the whole real cube.

The odd part of this atom is
\[
 \frac12\left(
   \frac{a_0+a\cdot t}{1+v\cdot t}
  -\frac{a_0-a\cdot t}{1-v\cdot t}\right)
 =\frac{(a-a_0v)\cdot t}{1-(v\cdot t)^2}.
\]
Set $u=a-a_0v$ and $w_i=|v_i|$. Uniform orientation gives
$(v\cdot t)^2=(w\cdot t)^2$, including the case $v=0$.
The arbitrary bias cancels when the full score is replaced by
$G(t)=(F(t)-F(-t))/2$.
If $f(t)=1$, then $F(t)>0$ and $F(-t)\leq0$, so $G(t)>0$.
If $f(t)=-1$, the reverse inequalities give $G(t)<0$.
This proves the forward statement directly under the model's original
weak inequality at negative labels.

For the converse, first suppose every $w_{hi}>0$. Each kernel has the
exact decomposition
\[
 \frac{u_h\cdot t}{1-(w_h\cdot t)^2}
 =\frac12\frac{u_h\cdot t}{1+w_h\cdot t}
  +\frac12\frac{u_h\cdot t}{1-w_h\cdot t}.
\]
In raw bits these two denominators are
\[
 1-\sum_iw_{hi}+2\sum_iw_{hi}x_i,
 \qquad
 1+\sum_iw_{hi}-2\sum_iw_{hi}x_i.
\]
Their slopes are respectively strictly positive and strictly negative;
both are positive throughout the cube by $\sum_iw_{hi}<1$.
Their numerators are arbitrary affine functions. The existing exact
positive/negative affine-ratio realization therefore implements the two
fractions by two original atoms.

For weakly positive $w_h$, replace each coordinate by $w_{hi}+\varepsilon$.
For sufficiently small positive $\varepsilon$, all sum bounds remain
strict. Every denominator stays positive and every score converges to
its original value as $\varepsilon\downarrow0$. There are only finitely
many vertices and all their signed scores are strictly positive, so
sufficiently small $\varepsilon$ preserves every sign. Apply the preceding
two-atom construction. This uses exactly two heads per original kernel
and does not change the input dimension.
\end{proof}

\paragraph{Interpretation as tangent kernels.}
For $w\in(0,\infty)^n$ with $\sum_iw_i<1$, differentiation gives
\[
 D_u\operatorname{atanh}(w\cdot t)
 =\frac{u\cdot t}{1-(w\cdot t)^2}.
\]
This interpretation is only a convenient identity; the proof above uses
no transcendental operation in the model.

\begin{proposition}[First Fourier level]
Let
\[
 K(w)_{ij}=2^{-n}\sum_{t\in\{-1,1\}^n}
       \frac{t_i t_j}{1-(w\cdot t)^2}.
\]
The degree-one Fourier coefficient vector of one kernel is $K(w)u$.
The real symmetric matrix $K(w)$ is positive definite and satisfies
$K(w)\succeq I$. Consequently, for fixed $w$, the degree-one coefficient
vector ranges over all of $\mathbb R^n$ as $u$ varies.
\end{proposition}

\begin{proof}
The coefficient identity follows by expanding $u\cdot t$ in the finite
sum. For any $z\in\mathbb R^n$,
\[
 z^TK(w)z
 =2^{-n}\sum_t\frac{(z\cdot t)^2}{1-(w\cdot t)^2}
 \geq 2^{-n}\sum_t(z\cdot t)^2=\sum_i z_i^2.
\]
Every denominator belongs to $(0,1]$, and the last equality uses
orthogonality of the coordinate characters. Thus $K(w)$ is invertible.
\end{proof}

For a sum of $H$ kernels, vanishing first-level Fourier coefficients is
exactly the vector equation $\sum_hK(w_h)u_h=0$. For fixed denominators
this is an $n$-equation constraint on the $nH$ numerator coefficients,
with full rank whenever $H\geq1$. It can be solved for the numerator
of one head. However, doing so does not generally preserve the target
signs, so it is not a permissible normalization of an arbitrary existing
small representation.

\begin{lemma}[Why first-level elimination is not automatic]
If an odd target is strictly sign-represented by a nonzero linear form
$\ell(t)=b\cdot t$, no strict score for that same target has all
degree-one Fourier coefficients equal to zero.
\end{lemma}

\begin{proof}
For every vertex $t$, strict signs imply $G(t)\ell(t)>0$.
Its average is therefore positive. But if all degree-one coefficients
of $G$ vanish, that average equals $\sum_i b_i\widehat G(\{i\})=0$,
a contradiction.
\end{proof}

\paragraph{An exact cancellation example.}
Conversely, low-degree cancellations can expose higher characters at
very small cost. On three bits, choose $w=(1/4,1/4,1/4)$ and $u=e_1$.
Direct calculation gives
\[
 \frac{t_1}{1-(t_1+t_2+t_3)^2/16}
 -\frac{48}{35}t_1-\frac{32}{105}(t_2+t_3)
 =\frac{32}{105}t_1t_2t_3
 \quad(t\in\{-1,1\}^3).
\]
The subtracted linear form is a second kernel with $w=0$.
Thus every degree-one coefficient cancels, while a strict parity score
remains. The converse theorem realizes its signs using four strict
original heads after an arbitrarily small perturbation. No positivity
of a kernel's Taylor coefficients excludes these signed cancellations.

\paragraph{Remaining lower-bound problem.}
An eventual lower bound of order $2^n/n$ for this positive-kernel class
on odd targets would imply the requested unrestricted lower bound,
since every original representation yields a representation in this
class with the same $H$. The converse shows that the two complexities
on odd targets differ by at most a factor of two. The normalization is
therefore global, not a restricted-head obstruction. A quantitative
bound on its strict sign patterns, uniform in all positive denominator
parameters and all numerator cancellations, is still missing.
\end{document}
